% Limitations paragraph for the threats-to-validity section: what the reported
% storage and energy figures are, and what they are not. No em dashes.
% Numbers verified against results/Journal_results/encoder_storage.json and the
% model_size records in results/*/seed_*/results.json.
\paragraph{Reported footprint and energy are analytical projections.}
Two accounting caveats bound the deployment claims in this paper. First, the
storage figures are projections from parameter and table counts rather than
measured serialized artifacts. Our int8 stage is a per-tensor fake-quantization
stand-in that rounds each linear layer's weights onto an 8-bit grid while
retaining floating-point tensors, so no checkpoint on disk is actually the size
we report; a production deployment would use quantization-aware training and
true 8-bit packing. The int8 column is therefore a lower bound in two respects:
layer biases are not quantized by our implementation, and the encoder's split
thresholds are retained at full precision because inference performs a
comparison search against them. Second, the deployed Tree-to-Spike artifact is
not the student network alone. The frozen encoder must ship with it, since
every incoming flow is binned against the teacher-derived split thresholds
before any spike is emitted, and the parameter-matched DNN-Tree twin consumes
the same encoder output and therefore carries the identical table cost. We
account for both components in Table~\ref{tab:architecture}. The tables are
not a rounding error: across the eleven datasets they add between $1.2$ and
$16.4$\,KB to the student weights, average $8.4$\,KB against the student's
$8.4$\,KB, and exceed the student on five of the eleven datasets, reaching
$67\%$ of the artifact on CICIoT2023. The complete deployment artifact
therefore averages approximately $17$\,KB (range $7$ to $30$\,KB) rather than
the $8$\,KB of student weights alone. Reporting only the network weights would
understate the deployed footprint by roughly a factor of two, which is the
same category of omission that makes spike counts a misleading efficiency
proxy for hybrid architectures.
Finally, as stated in Section~\ref{sec:energy}, the operation counts behind the
energy figures are measured on the test sets, but they are priced with a
published per-operation model rather than measured on neuromorphic silicon;
absolute joule claims would require hardware measurement on a device such as
Loihi or Akida.
